\medskip

\textbf{Partie A}

\medskip

On considère la fonction $f$ définie pour tout réel $x$ de ]0~;~1] par:

\[f(x) = \text{e}^{-x} + \ln (x).\]

\smallskip

\begin{enumerate}
\item Calculer la limite de $f$ en $0$.
\item On admet que $f$ est dérivable sur ]0~;~1]. On note $f'$ sa fonction dérivée.

Démontrer que, pour tout réel $x$ appartenant à ]0~;~1], on a :

\[f'(x) = \dfrac{1 - x\text{e}^{-x}}{x}\]

\item Justifier que, pour tout réel $x$ appartenant à ]0~;~1], on a $x\text{e}^{-x} < 1$.

En déduire le tableau de variation de $f$ sur ]0~;~1].
\item Démontrer qu'il existe un unique réel $\ell$ appartenant à ]0~;~1] tel que $f(\ell) = 0$.
\end{enumerate}

\medskip

\textbf{Partie B}

\medskip

\begin{enumerate}
\item On définit deux suites $\left(a_n\right)$ et $\left(b_n\right)$ par:

\[\left\{\begin{array}{l c l}
a_0& =& \dfrac{1}{10}\\
b_0& =& 1
\end{array}\right.\: \text{et, pour tout entier naturel }\: n,\left\{\begin{array}{l c l}
a_{n+1}&=&\text{e}^{-b_n}\\
b_{n+1}&=&\text{e}^{-a_n}
\end{array}\right.\]

	\begin{enumerate}
		\item Calculer $a_1$ et $b_1$. On donnera des valeurs approchées à $10^{-2}$ près.
		\item On considère ci-dessous la fonction \texttt{termes}, écrite en langage Python.
		
\begin{center}

\begin{tabular}{|l|}\hline
\texttt{def termes (n) :}\\
\quad \texttt{a=1/10}\\
\quad \texttt{b=1}\\
\quad \texttt{for k in range(0,n) :}\\
\qquad \texttt{c= \ldots}\\
\qquad \texttt{b  = \ldots}\\
\qquad \texttt{a = c}\\
\quad \texttt{return(a,b)}\\ \hline
\end{tabular}
\end{center}

Recopier et compléter sans justifier le cadre ci-dessus de telle sorte que la fonction termes calcule les termes des suites $\left(a_n\right)$ et $\left(b_n\right)$.
	\end{enumerate}
\item On rappelle que la fonction $x \longmapsto \text{e}^{-x}$ est décroissante sur $\R$.
	\begin{enumerate}
		\item Démontrer par récurrence que, pour tout entier naturel $n$, on a :
		
		\[0 < a_n \leqslant a_{n+1} \leqslant b_{n+1} \leqslant b_n \leqslant 1\]
		
		\item En déduire que les suites $\left(a_n\right)$ et $\left(b_n\right)$ sont convergentes.
	\end{enumerate}	
\item On note $A$ la limite de $\left(a_n\right)$ et $B$ la limite de $\left(b_n\right)$.

On admet que $A$ et $B$ appartiennent à l'intervalle ]0~;~1], et que $A = \text{e}^{-B}$ et $B = \text{e}^{-A}$.
	\begin{enumerate}
		\item Démontrer que $f(A) = 0$.
		\item Déterminer $A - B$.
	\end{enumerate}
\end{enumerate}
